% ============================================================
\chapter{Financial Markets --- Modelling}
\label{ch:markets}
% ============================================================

In 1900, Louis Bachelier defended his thesis at the Sorbonne under Henri
Poincar\'e's supervision. His subject: the ``theory of speculation,'' in
which he proposed to model stock prices as a stochastic process. Poincar\'e,
though visionary, judged the work ``not without interest'' but awarded it an
honorable rather than a very honorable mention. It took sixty years, and
the arrival of Samuelson, Black, Scholes, and Merton, for Bachelier's ideas
to become the cornerstone of modern finance. Today, financial markets move
trillions each day, and the mathematical models that describe them lie at
the heart of the global economy.

% ------------------------------------------------------------
\section{Introduction to Financial Markets}
% ------------------------------------------------------------

Financial markets are venues for the exchange of assets whose value is
uncertain. Understanding their functioning and modelling them mathematically
is the starting point of all quantitative finance.

\begin{definition}[Financial Market]
A \textbf{financial market} is an organised system enabling the exchange
of financial assets between economic agents. We distinguish:
\begin{itemize}
    \item The \textbf{primary market}: issuance of new securities.
    \item The \textbf{secondary market}: trading of previously issued securities.
\end{itemize}
\end{definition}

\subsection{Asset Classes}

\begin{definition}[Financial Asset]
A \textbf{financial asset} is a contract entitling its holder to future
cash flows. The main classes are:
\begin{enumerate}
    \item \textbf{Equities}: ownership shares in a company.
    \item \textbf{Bonds}: fixed-income debt securities.
    \item \textbf{Derivatives}: contracts whose value depends on an underlying asset.
    \item \textbf{Currencies}: foreign exchange (forex) market.
    \item \textbf{Commodities}: gold, oil, wheat, etc.
\end{enumerate}
\end{definition}

\subsection{Fundamental Derivatives}

\begin{definition}[European Option]
A \textbf{European call option} with strike price $K$ and maturity $T$ on
an underlying $S$ gives its holder the right (but not the obligation) to buy
$S$ at price $K$ at time $T$. Its payoff is:
\[
    (S_T - K)^+ = \max(S_T - K, 0).
\]
Similarly, a \textbf{European put option} has payoff:
\[
    (K - S_T)^+ = \max(K - S_T, 0).
\]
\end{definition}

\begin{definition}[Forward Contract]
A \textbf{forward contract} is an agreement to buy or sell an asset at
a future date $T$ at a price $F$ fixed today. The buyer's payoff is $S_T - F$.
\end{definition}

\section{Discrete-Time Modelling}

\subsection{One-Period Model}

Consider a market with two dates $t=0$ and $t=1$, with:
\begin{itemize}
    \item A risk-free asset with price $B_0 = 1$, $B_1 = 1+r$.
    \item A risky asset with price $S_0$ at $t=0$, which can take values
          $S_1^u = uS_0$ (up) or $S_1^d = dS_0$ (down) with $0 < d < 1+r < u$.
\end{itemize}

\begin{definition}[Trading Strategy]
A \textbf{trading strategy} is a pair $(\phi^0, \phi^1) \in \R^2$ where
$\phi^0$ represents the quantity of risk-free asset and $\phi^1$ the quantity
of risky asset held. The portfolio value is:
\[
    V_t = \phi^0 B_t + \phi^1 S_t, \quad t \in \{0,1\}.
\]
\end{definition}

\begin{definition}[Arbitrage Opportunity]
An \textbf{arbitrage opportunity} is a strategy $(\phi^0, \phi^1)$ such that:
\[
    V_0 \leq 0, \quad V_1 \geq 0 \text{ a.s.}, \quad \PP(V_1 > 0) > 0.
\]
That is, a possible gain without risk of loss and without initial investment.
\end{definition}

\begin{theorem}[No-Arbitrage --- One-Period Binomial Model]
The one-period binomial model is arbitrage-free if and only if $d < 1+r < u$.
\end{theorem}

\begin{proof}
Suppose there exists an arbitrage $(\phi^0, \phi^1)$ with $V_0 = 0$. Then:
\[
    \phi^0 + \phi^1 S_0 = 0 \implies \phi^0 = -\phi^1 S_0.
\]
The terminal values are:
\begin{align*}
    V_1^u &= \phi^1(uS_0 - (1+r)S_0) = \phi^1 S_0(u - 1 - r), \\
    V_1^d &= \phi^1(dS_0 - (1+r)S_0) = \phi^1 S_0(d - 1 - r).
\end{align*}
For both $V_1^u \geq 0$ and $V_1^d \geq 0$ with at least one strict
inequality, $u - 1 - r$ and $d - 1 - r$ must have the same non-zero sign,
contradicting $d < 1+r < u$.
\end{proof}

\subsection{Risk-Neutral Probability}

\begin{theorem}[Risk-Neutral Measure]
\label{thm:risk_neutral}
Under the no-arbitrage condition ($d < 1+r < u$), there exists a unique
probability measure $\mathbb{Q}$ on $\{\omega_u, \omega_d\}$ such that the
discounted risky asset price is a $\mathbb{Q}$-martingale:
\[
    S_0 = \frac{1}{1+r}\E^\mathbb{Q}[S_1].
\]
The risk-neutral probability is given by:
\[
    q = \frac{(1+r) - d}{u - d}, \quad 1-q = \frac{u - (1+r)}{u - d}.
\]
\end{theorem}

\begin{keyformulas}[title=\textbf{Risk-Neutral Pricing of a Contingent Claim}]
The price at $t=0$ of a contingent claim with payoff $H = h(S_1)$ is:
\[
    V_0 = \frac{1}{1+r}\E^\mathbb{Q}[H]
        = \frac{1}{1+r}\bigl[q \cdot h(uS_0) + (1-q) \cdot h(dS_0)\bigr].
\]
\end{keyformulas}

\begin{example}
Let $S_0 = 100$, $u = 1.2$, $d = 0.9$, $r = 0.05$, $K = 100$.

The risk-neutral probability is:
\[
    q = \frac{1.05 - 0.9}{1.2 - 0.9} = \frac{0.15}{0.30} = 0.5.
\]
The call price is:
\[
    C_0 = \frac{1}{1.05}\bigl[0.5 \times (120-100)^+ + 0.5 \times (90-100)^+\bigr]
        = \frac{0.5 \times 20}{1.05} \approx 9.52.
\]
\end{example}

\section{Multi-Period Binomial Model (CRR)}

\begin{definition}[Cox--Ross--Rubinstein Model]
The \textbf{CRR model} with $N$ periods is defined by:
\begin{itemize}
    \item $S_{n+1} = S_n \cdot Y_{n+1}$ where $Y_{n+1} \in \{u, d\}$
          are i.i.d.\ under $\mathbb{Q}$.
    \item The discounted price $\tilde{S}_n = S_n / (1+r)^n$ is a
          $\mathbb{Q}$-martingale.
    \item After $n$ periods: $S_n = S_0 u^j d^{n-j}$ for $j$ up-moves
          out of $n$.
\end{itemize}
\end{definition}

\begin{theorem}[CRR Pricing Formula]
The price of a European call in the CRR model is:
\[
    C_0 = \frac{1}{(1+r)^N} \sum_{j=0}^{N} \binom{N}{j} q^j (1-q)^{N-j}
    \bigl(S_0 u^j d^{N-j} - K\bigr)^+.
\]
\end{theorem}

\begin{intuition}[title=\textbf{Convergence to Black--Scholes}]
As $N \to \infty$ with the parametrisation $u = e^{\sigma\sqrt{\Delta t}}$,
$d = e^{-\sigma\sqrt{\Delta t}}$, $\Delta t = T/N$, the CRR formula converges
to the Black--Scholes formula (Chapter 2). This is a consequence of the
Central Limit Theorem.
\end{intuition}

\subsection{CRR Parameter Calibration}

For the discrete model to approximate the continuous model with volatility
$\sigma$ and rate $r$, we set:
\begin{align}
    u &= e^{\sigma\sqrt{\Delta t}}, \quad d = e^{-\sigma\sqrt{\Delta t}} = 1/u, \\
    q &= \frac{e^{r\Delta t} - d}{u - d}.
\end{align}

\section{Fundamental Theorems of Asset Pricing}

\begin{theorem}[First Fundamental Theorem of Asset Pricing (FTAP)]
\label{thm:ftap1}
A market is arbitrage-free if and only if there exists a probability measure
$\mathbb{Q}$ equivalent to $\PP$ under which discounted asset prices are
martingales.
\end{theorem}

\begin{theorem}[Second Fundamental Theorem]
\label{thm:ftap2}
An arbitrage-free market is \textbf{complete} (every contingent claim is
replicable) if and only if the martingale measure $\mathbb{Q}$ is unique.
\end{theorem}

\begin{remark}
In the binomial model, the market is complete because there are two states
of nature and two assets. In contrast, in a trinomial model (three states),
the market is incomplete and infinitely many martingale measures exist.
\end{remark}

\section{Replication and Hedging}

\begin{proposition}[Replicating Portfolio --- One Period]
In the one-period binomial model, the contingent claim $H = h(S_1)$ is
replicated by $(\phi^0, \phi^1)$ with:
\[
    \phi^1 = \frac{h(uS_0) - h(dS_0)}{(u-d)S_0}, \quad
    \phi^0 = \frac{1}{1+r}\bigl[h(uS_0) - \phi^1 uS_0\bigr].
\]
The quantity $\phi^1$ is called the \textbf{delta} of the option.
\end{proposition}

\section{Python Implementation}

\begin{codeblock}[title=\textbf{CRR Binomial Model in Python}]
\begin{minted}{python}
import numpy as np

def crr_european(S0, K, T, r, sigma, N, option='call'):
    """Price a European option using the CRR binomial model."""
    dt = T / N
    u = np.exp(sigma * np.sqrt(dt))
    d = 1 / u
    q = (np.exp(r * dt) - d) / (u - d)
    discount = np.exp(-r * dt)

    # Terminal prices
    j = np.arange(N + 1)
    ST = S0 * u**j * d**(N - j)

    if option == 'call':
        payoff = np.maximum(ST - K, 0)
    else:
        payoff = np.maximum(K - ST, 0)

    # Backward induction
    for n in range(N - 1, -1, -1):
        payoff = discount * (q * payoff[1:] + (1 - q) * payoff[:-1])

    return payoff[0]

# Example
S0, K, T, r, sigma = 100, 100, 1.0, 0.05, 0.20
for N in [10, 50, 100, 500, 1000]:
    price = crr_european(S0, K, T, r, sigma, N)
    print(f"N={N:4d}: C = {price:.4f}")
\end{minted}
\end{codeblock}

\begin{codeblock}[title=\textbf{Put-Call Parity --- Verification}]
\begin{minted}{python}
def verify_put_call_parity(S0, K, T, r, sigma, N):
    """Verify put-call parity: C - P = S0 - K*exp(-rT)."""
    C = crr_european(S0, K, T, r, sigma, N, 'call')
    P = crr_european(S0, K, T, r, sigma, N, 'put')
    parity = S0 - K * np.exp(-r * T)
    print(f"C - P = {C - P:.6f}")
    print(f"S0 - K*exp(-rT) = {parity:.6f}")
    print(f"Error = {abs(C - P - parity):.2e}")

verify_put_call_parity(100, 100, 1.0, 0.05, 0.20, 500)
\end{minted}
\end{codeblock}

\section{Put-Call Parity}

\begin{theorem}[Put-Call Parity]
For European options with the same strike $K$ and maturity $T$:
\[
    C - P = S_0 - K e^{-rT}.
\]
\end{theorem}

\begin{proof}
Consider the portfolio $\Pi = C - P - S_0 + Ke^{-rT}$. At maturity:
\[
    \Pi_T = (S_T - K)^+ - (K - S_T)^+ - S_T + K = S_T - K - S_T + K = 0.
\]
By no-arbitrage, $\Pi_0 = 0$, yielding the result.
\end{proof}

\section{Market Hypotheses}

\begin{warning}[title=\textbf{Simplifying Assumptions}]
Classical models rely on strong assumptions:
\begin{enumerate}
    \item Frictionless market (no transaction costs).
    \item Continuous trading is possible.
    \item Short selling is allowed without restrictions.
    \item Borrowing rate equals lending rate.
    \item No counterparty risk.
\end{enumerate}
Relaxing these assumptions leads to more realistic but significantly
more complex models.
\end{warning}

\section{Stylised Facts of Financial Returns}

Log-returns $r_t = \ln(S_t/S_{t-1})$ exhibit well-documented empirical
properties:

\begin{enumerate}
    \item \textbf{Heavy tails}: the return distribution has heavier tails
          than the normal distribution (excess kurtosis).
    \item \textbf{Asymmetry}: slight negative skewness.
    \item \textbf{Aggregational Gaussianity}: returns over longer horizons
          approach normality.
    \item \textbf{Absence of autocorrelation}: returns are nearly uncorrelated.
    \item \textbf{Volatility clustering}: periods of high volatility are
          clustered together.
    \item \textbf{Leverage effect}: negative correlation between returns
          and volatility.
\end{enumerate}

\begin{codeblock}[title=\textbf{Stylised Facts Analysis}]
\begin{minted}{python}
import numpy as np
from scipy import stats

def stylized_facts(returns):
    """Analyse stylised facts of returns."""
    print(f"Mean:       {np.mean(returns):.6f}")
    print(f"Std dev:    {np.std(returns):.6f}")
    print(f"Skewness:   {stats.skew(returns):.4f}")
    print(f"Kurtosis:   {stats.kurtosis(returns):.4f}")

    # Normality test (Jarque-Bera)
    jb_stat, jb_pval = stats.jarque_bera(returns)
    print(f"Jarque-Bera: stat={jb_stat:.2f}, p={jb_pval:.4e}")

    # Autocorrelation of returns and squared returns
    n = len(returns)
    r_centered = returns - np.mean(returns)
    acf1 = np.correlate(r_centered[:-1], r_centered[1:]) / (n * np.var(returns))
    r2 = r_centered**2
    acf1_sq = np.correlate(r2[:-1], r2[1:]) / (n * np.var(r2))
    print(f"ACF(1) returns:  {acf1[0]:.4f}")
    print(f"ACF(1) squared:  {acf1_sq[0]:.4f}")

# Simulated GBM for comparison
np.random.seed(42)
S0, mu, sigma, dt, N = 100, 0.08, 0.20, 1/252, 2520
Z = np.random.standard_normal(N)
log_returns = (mu - 0.5*sigma**2)*dt + sigma*np.sqrt(dt)*Z
stylized_facts(log_returns)
\end{minted}
\end{codeblock}

\section{Exercises}

\begin{exercise}[One-Period Binomial Model]
Let $S_0 = 50$, $u = 1.3$, $d = 0.8$, $r = 0.04$.
\begin{enumerate}
    \item Verify the no-arbitrage condition.
    \item Compute the risk-neutral probability $q$.
    \item Compute the price of a call with strike $K = 55$.
    \item Determine the replicating portfolio $(\phi^0, \phi^1)$.
    \item Verify that the put price satisfies put-call parity.
\end{enumerate}
\end{exercise}

\begin{exercise}[CRR Convergence]
\begin{enumerate}
    \item Implement the CRR model and plot the call price as a function
          of $N$ for $N = 1, 2, \ldots, 200$.
    \item Observe the convergence and even/odd oscillations.
    \item Compare with the Black--Scholes formula (Chapter 2).
\end{enumerate}
\end{exercise}

\begin{exercise}[Incomplete Market]
Consider a trinomial model: $S_1 \in \{uS_0, mS_0, dS_0\}$ with $d < m < u$.
\begin{enumerate}
    \item Show that the market is incomplete.
    \item Characterise the set of martingale measures.
    \item Deduce the no-arbitrage bounds for the call price.
\end{enumerate}
\end{exercise}

\begin{exercise}[American Options in the CRR Model]
\begin{enumerate}
    \item Modify the CRR algorithm to price an American put by incorporating
          the possibility of early exercise.
    \item Compare American and European put prices. Verify that the American
          put is always at least as expensive.
    \item Determine the optimal exercise boundary.
\end{enumerate}
\end{exercise}
